\documentclass[a4paper,12pt]{article}

\usepackage[T2A]{fontenc}
\usepackage[utf8]{inputenc}
\usepackage[russian]{babel}

\usepackage{amsmath,amssymb,amsthm,mathtools}
\usepackage{helvet}
\renewcommand{\familydefault}{\sfdefault}
\usepackage{icomma}
\usepackage{geometry}
\usepackage{microtype}
\usepackage{tikz}
\usetikzlibrary{arrows.meta,calc}
\usepackage{tkz-euclide}
\usepackage{pgfplots}
\pgfplotsset{compat=1.18}
\usepackage{tabularx,booktabs,longtable}
\usepackage{enumitem}
\usepackage{hyperref}
\usepackage{parskip}
\usepackage{fancyhdr}
\usepackage{setspace}
\usepackage{xcolor}
\usepackage{titlesec}

\geometry{left=18mm,right=18mm,top=18mm,bottom=20mm}

\definecolor{accent}{HTML}{008080}
\definecolor{darkaccent}{HTML}{004D4D}
\definecolor{textgray}{HTML}{333333}
\definecolor{softaccent}{HTML}{EAF5F5}

\hypersetup{
    colorlinks=true,
    linkcolor=darkaccent,
    urlcolor=darkaccent
}

\onehalfspacing
\setlength{\parindent}{0pt}
\setlength{\parskip}{6pt}

\pagestyle{fancy}
\fancyhf{}
\lhead{\textcolor{textgray}{Чек-лист}}
\rhead{\textcolor{textgray}{ЕГЭ: тригонометрические уравнения}}
\cfoot{\textcolor{textgray}{\thepage}}

\titleformat{\section}
  {\Large\bfseries\color{darkaccent}}
  {\thesection}{0.7em}{}
  [\vspace{-1mm}{\color{accent}\titlerule[0.6pt]}]

\titleformat{\subsection}
  {\large\bfseries\color{darkaccent}}
  {\thesubsection}{0.7em}{}

\setlist[itemize]{leftmargin=7mm,itemsep=2mm,topsep=2mm}
\setlist[enumerate]{leftmargin=8mm,itemsep=2mm,topsep=2mm}

\providecommand{\tg}{\mathop{\mathrm{tg}}\nolimits}
\providecommand{\ctg}{\mathop{\mathrm{ctg}}\nolimits}
\DeclareMathOperator{\arctg}{arctg}

\newcommand{\important}[1]{\textcolor{darkaccent}{\textbf{#1}}}

\begin{document}

\begin{titlepage}
\thispagestyle{empty}
\begin{center}
\vspace*{22mm}

{\Huge\bfseries\color{darkaccent} Чек-лист}

\vspace{7mm}

{\Large\bfseries
Тригонометрические уравнения:\\
разложение на множители, группировка и ограничения
}

\vspace{18mm}

{\large
Лёвин Артём Александрович\\
эксклюзивно для Тимофея Васильченко
}

\vspace{8mm}

{\large \today}

\vfill

\begin{tikzpicture}[x=1cm,y=1cm]
  \draw[accent,line width=1.1pt]
    (-5,0) .. controls (-3.4,0.9) and (-1.8,-0.7) .. (-0.2,0.2)
           .. controls (1.5,1.2) and (3.1,-0.8) .. (5,0.15);
  \draw[darkaccent,line width=0.45pt]
    (-4.5,-0.35) .. controls (-2.8,0.25) and (-1.1,-0.55) .. (0.5,-0.2)
           .. controls (2.1,0.2) and (3.4,-0.45) .. (4.5,-0.15);
\end{tikzpicture}

\vspace{12mm}
\end{center}
\end{titlepage}

\section{Главная идея}

В тригонометрических уравнениях сначала нужно добиться двух вещей:

\begin{enumerate}
  \item \important{одинаковых углов};
  \item \important{одинаковых функций или удобного разложения на множители}.
\end{enumerate}

Если уравнение содержит $\sin 2x$, $\cos 2x$, $\sin^2x$, $\cos^2x$, сдвиги вида $\frac{\pi}{2}+x$, $\frac{3\pi}{2}+x$, то почти всегда нужно применить формулы преобразования.

\important{Главный принцип:} не пытайтесь решать уравнение ``как есть''. Сначала приведите его к виду, где можно вынести общий множитель, сгруппировать слагаемые или сделать замену.

\section{Одинаковые углы}

\subsection{Формула двойного угла}

Если в уравнении есть углы $x$ и $2x$, обычно переходят от $2x$ к $x$:

\[
\sin 2x=2\sin x\cos x.
\]

Пример:
\[
\sin 2x+\sqrt3\sin x=0.
\]

Преобразуем:
\[
2\sin x\cos x+\sqrt3\sin x=0,
\]
\[
\sin x(2\cos x+\sqrt3)=0.
\]

Получаем совокупность:
\[
\begin{cases}
\sin x=0,\\
2\cos x+\sqrt3=0.
\end{cases}
\]

Отсюда:
\[
x=\pi n,\quad n\in\mathbb Z,
\]
или
\[
\cos x=-\frac{\sqrt3}{2}.
\]

Тогда
\[
x=\frac{5\pi}{6}+2\pi k
\quad\text{или}\quad
x=\frac{7\pi}{6}+2\pi k,
\qquad k\in\mathbb Z.
\]

\subsection{Что важно}

После перехода
\[
\sin 2x=2\sin x\cos x
\]
не нужно раскрывать всё дальше. Наоборот, нужно искать общий множитель.

\important{Цель:} получить произведение, равное нулю.

\section{Группировка и двойное вынесение}

\subsection{Когда работает группировка}

Если после переноса всех слагаемых в одну сторону появилось четыре слагаемых, часто помогает группировка:

\[
A+B+C+D=0.
\]

Обычно группируют так, чтобы внутри двух скобок появился одинаковый множитель.

\subsection{Типовой пример}

Решим уравнение:
\[
\sin 2x+\sqrt2\sin x=2\cos x+\sqrt2.
\]

Переносим всё влево:
\[
\sin 2x+\sqrt2\sin x-2\cos x-\sqrt2=0.
\]

Раскрываем $\sin 2x$:
\[
2\sin x\cos x+\sqrt2\sin x-2\cos x-\sqrt2=0.
\]

Группируем:
\[
(2\sin x\cos x+\sqrt2\sin x)+(-2\cos x-\sqrt2)=0.
\]

Выносим:
\[
\sin x(2\cos x+\sqrt2)-(2\cos x+\sqrt2)=0.
\]

Выносим общий множитель:
\[
(2\cos x+\sqrt2)(\sin x-1)=0.
\]

Получаем:
\[
2\cos x+\sqrt2=0
\quad\text{или}\quad
\sin x-1=0.
\]

Ответ:
\[
x=\frac{3\pi}{4}+2\pi k,\quad
x=\frac{5\pi}{4}+2\pi k,\quad
x=\frac{\pi}{2}+2\pi k,
\qquad k\in\mathbb Z.
\]

\subsection{Техническое правило}

Если группируете так:
\[
A+B-C-D,
\]
то удобно писать:
\[
(A+B)+(-C-D).
\]

Минус нельзя ``отрывать'' от слагаемых. Иначе легко получить ошибку знака.

\section{Формулы суммы и разности}

\subsection{Базовые формулы}

\[
\sin(\alpha+\beta)=\sin\alpha\cos\beta+\cos\alpha\sin\beta.
\]

\[
\cos(\alpha+\beta)=\cos\alpha\cos\beta-\sin\alpha\sin\beta.
\]

\[
\cos(\alpha-\beta)=\cos\alpha\cos\beta+\sin\alpha\sin\beta.
\]

\subsection{Частые преобразования}

\[
\sin\left(\frac{3\pi}{2}+x\right)=-\cos x.
\]

\[
\cos\left(\frac{\pi}{2}+x\right)=-\sin x.
\]

\[
\cos(2\pi-x)=\cos x.
\]

Эти равенства можно получать через формулы суммы и разности или через единичную окружность.

\subsection{Мини-пример}

Упростим выражение:
\[
2\sin x\sin\left(\frac{3\pi}{2}+x\right)
-\sqrt3\cos x\cos\left(\frac{\pi}{2}+x\right).
\]

Используем:
\[
\sin\left(\frac{3\pi}{2}+x\right)=-\cos x,
\qquad
\cos\left(\frac{\pi}{2}+x\right)=-\sin x.
\]

Тогда:
\[
2\sin x(-\cos x)-\sqrt3\cos x(-\sin x)
=
-\!2\sin x\cos x+\sqrt3\sin x\cos x.
\]

Получаем:
\[
(\sqrt3-2)\sin x\cos x.
\]

\section{Одинаковая функция и замена}

\subsection{Когда нужна замена}

Если углы уже одинаковые, но функции разные, нужно подумать, можно ли выразить одну функцию через другую.

Главная формула:
\[
\sin^2x+\cos^2x=1.
\]

Отсюда:
\[
\cos^2x=1-\sin^2x,
\qquad
\sin^2x=1-\cos^2x.
\]

\subsection{Пример}

Рассмотрим уравнение:
\[
2\sin^3x-\sin^2x+2\sin x\cos^2x-\cos^2x=0.
\]

Заменим:
\[
\cos^2x=1-\sin^2x.
\]

Получаем:
\[
2\sin^3x-\sin^2x+2\sin x(1-\sin^2x)-(1-\sin^2x)=0.
\]

Раскрываем:
\[
2\sin^3x-\sin^2x+2\sin x-2\sin^3x-1+\sin^2x=0.
\]

Сокращаем:
\[
2\sin x-1=0.
\]

Отсюда:
\[
\sin x=\frac12.
\]

Ответ:
\[
x=\frac{\pi}{6}+2\pi k
\quad\text{или}\quad
x=\frac{5\pi}{6}+2\pi k,
\qquad k\in\mathbb Z.
\]

\subsection{Если после замены получается многочлен}

Если Вы обозначили
\[
t=\sin x
\]
или
\[
t=\cos x,
\]
то после решения уравнения относительно $t$ обязательно возвращайтесь к тригонометрии:

\[
t=1\quad\Rightarrow\quad \sin x=1.
\]

\[
t=-\frac12\quad\Rightarrow\quad \cos x=-\frac12.
\]

\important{Нельзя оставить ответ только в виде значений $t$.}

\section{Ограничения в тригонометрических уравнениях}

\subsection{Когда появляются ограничения}

Ограничения обязательно проверяются, если есть:

\begin{itemize}
  \item корень чётной степени;
  \item знаменатель;
  \item $\tg x$ или $\ctg x$;
  \item логарифм;
  \item дробь, в знаменателе которой стоит логарифм или тригонометрическое выражение.
\end{itemize}

\subsection{Корень}

Если есть
\[
\sqrt{f(x)},
\]
то
\[
f(x)\ge0.
\]

Если корень стоит множителем, например
\[
\sqrt{\cos x}\,(\tg x+\sqrt3)(\cos x-1)=0,
\]
то сначала пишем ограничения:
\[
\cos x\ge0.
\]

Но из-за $\tg x$ дополнительно:
\[
\cos x\ne0.
\]

Значит, вместе:
\[
\cos x>0.
\]

\subsection{Тангенс}

\[
\tg x=\frac{\sin x}{\cos x},
\]
поэтому:
\[
\cos x\ne0.
\]

\[
\ctg x=\frac{\cos x}{\sin x},
\]
поэтому:
\[
\sin x\ne0.
\]

\subsection{Логарифм}

Для логарифма
\[
\log_a b
\]
должно быть:
\[
a>0,\qquad a\ne1,\qquad b>0.
\]

Если логарифм стоит в знаменателе, то дополнительно:
\[
\log_a b\ne0.
\]

А это означает:
\[
b\ne1.
\]

\section{Уравнения с ограничениями}

\subsection{Пример с корнем и тангенсом}

Решим:
\[
\sqrt{\cos x}\,(\tg x+\sqrt3)(\cos x-1)=0.
\]

Ограничения:
\[
\cos x\ge0,
\qquad
\cos x\ne0.
\]

Значит:
\[
\cos x>0.
\]

Теперь решаем уравнение по множителям:

\[
\cos x=1
\quad\Rightarrow\quad
x=2\pi k.
\]

\[
\tg x=-\sqrt3
\quad\Rightarrow\quad
x=-\frac{\pi}{3}+\pi k.
\]

\[
\cos x-1=0
\quad\Rightarrow\quad
x=2\pi k.
\]

Проверяем ограничение $\cos x>0$.

Для
\[
x=2\pi k
\]
имеем
\[
\cos x=1>0,
\]
значит, корни подходят.

Для
\[
x=-\frac{\pi}{3}+\pi k
\]
подходит только точка из правой полуплоскости:
\[
x=-\frac{\pi}{3}+2\pi k.
\]

Ответ:
\[
x=2\pi k
\quad\text{или}\quad
x=-\frac{\pi}{3}+2\pi k,
\qquad k\in\mathbb Z.
\]

\subsection{Отбор на окружности}

Условие
\[
\cos x>0
\]
означает правую полуплоскость единичной окружности:

\begin{center}
\begin{tikzpicture}[scale=1.0]
  \draw[-{Latex[length=2mm]}] (-2.2,0)--(2.2,0);
  \draw[-{Latex[length=2mm]}] (0,-2.2)--(0,2.2);
  \draw[accent,line width=0.8pt] (0,0) circle (1.6);
  \draw[darkaccent,line width=1.2pt] (0,1.6) arc[start angle=90,end angle=-90,radius=1.6];
  \node[right] at (2.2,0) {$\cos x>0$};
  \fill[darkaccent] (0.8,-1.3856) circle (2pt);
  \node[below right] at (0.8,-1.3856) {\small $-\frac{\pi}{3}$};
  \fill[darkaccent] (1.6,0) circle (2pt);
  \node[above right] at (1.6,0) {\small $0$};
  \node[left] at (-1.7,0) {$\pi$};
  \node[above] at (0,1.7) {$\frac{\pi}{2}$};
  \node[below] at (0,-1.7) {$-\frac{\pi}{2}$};
\end{tikzpicture}
\end{center}

\subsection{Пример с логарифмом в знаменателе}

Решим:
\[
\frac{(\tg x+\sqrt3)(2\sin^2x-2)}
{\log_{31}(\sqrt2\cos x)}=0.
\]

Сначала ограничения.

Тело логарифма положительно:
\[
\sqrt2\cos x>0
\quad\Rightarrow\quad
\cos x>0.
\]

Знаменатель не равен нулю:
\[
\log_{31}(\sqrt2\cos x)\ne0.
\]

Значит:
\[
\sqrt2\cos x\ne1,
\qquad
\cos x\ne\frac{\sqrt2}{2}.
\]

Также из-за $\tg x$:
\[
\cos x\ne0,
\]
но это уже учтено условием $\cos x>0$.

Теперь решаем числитель:
\[
(\tg x+\sqrt3)(2\sin^2x-2)=0.
\]

Первый множитель:
\[
\tg x=-\sqrt3
\quad\Rightarrow\quad
x=-\frac{\pi}{3}+\pi k.
\]

С учётом $\cos x>0$:
\[
x=-\frac{\pi}{3}+2\pi k.
\]

Второй множитель:
\[
2\sin^2x-2=0,
\]
\[
\sin^2x=1.
\]

Через понижение степени:
\[
\frac{1-\cos2x}{2}=1,
\]
\[
\cos2x=-1.
\]

Отсюда:
\[
2x=\pi+2\pi k,
\qquad
x=\frac{\pi}{2}+\pi k.
\]

Но для этих корней
\[
\cos x=0,
\]
поэтому они не подходят.

Ответ:
\[
x=-\frac{\pi}{3}+2\pi k,\qquad k\in\mathbb Z.
\]

\section{Типичные ошибки}

\begin{longtable}{@{}p{0.33\linewidth}p{0.61\linewidth}@{}}
\toprule
\textbf{Ошибка} & \textbf{Как правильно}\\
\midrule
Решать уравнение с разными углами без преобразований &
Сначала привести углы к одному виду: чаще всего через $\sin2x=2\sin x\cos x$.\\

Не выносить общий множитель &
После преобразования ищите множители: $\sin x$, $\cos x$, общую скобку.\\

Ошибиться при группировке &
Минус перед группой слагаемых переносится внутрь скобок: $-2\cos x-\sqrt2=-(2\cos x+\sqrt2)$.\\

Забыть ограничения из-за $\tg x$ &
Для $\tg x$ всегда требуется $\cos x\ne0$.\\

Забыть ограничения из-за корня &
Для $\sqrt{f(x)}$ обязательно $f(x)\ge0$.\\

Забыть, что логарифм в знаменателе не равен нулю &
Если $\log_a b$ стоит в знаменателе, нужно не только $b>0$, но и $\log_a b\ne0$, то есть $b\ne1$.\\

Не проверить найденные корни по ограничениям &
Сначала найдите все кандидаты, затем отбросьте те, которые нарушают ОДЗ.\\

Оставить лишнюю точку при отборе по окружности &
Если после ограничения подходит только одна из двух диаметрально противоположных точек, период меняется с $\pi$ на $2\pi$.\\
\bottomrule
\end{longtable}

\newpage

\section*{Тренировочный блок}

\textbf{Инструкция.} Решите уравнения. В заданиях с ограничениями сначала выпишите ОДЗ.

\subsection*{Средний уровень}

\begin{enumerate}
  \item Решите уравнение:
  \[
  \sin2x+\sqrt3\sin x=0.
  \]

  \item Решите уравнение:
  \[
  \sin2x+\sqrt2\sin x=2\cos x+\sqrt2.
  \]

  \item Решите уравнение:
  \[
  2\cos^2x-\sqrt3\cos x=0.
  \]

\subsection*{Выше среднего уровня}

  \item Решите уравнение:
  \[
  \sin2x-\sin x-2\cos x+1=0.
  \]

  \item Решите уравнение:
  \[
  2\sin^3x-\sin^2x+2\sin x\cos^2x-\cos^2x=0.
  \]

  \item Решите уравнение:
  \[
  \sin\left(\frac{3\pi}{2}+x\right)+\cos\left(\frac{\pi}{2}+x\right)=0.
  \]

  \item Решите уравнение:
  \[
  \sqrt{\cos x}\,(\tg x+\sqrt3)(\cos x-1)=0.
  \]

  \item Решите уравнение:
  \[
  \sqrt{\sin x}\,(2\cos x-1)(\sin x-1)=0.
  \]

  \item Решите уравнение:
  \[
  \frac{(\tg x+\sqrt3)(2\sin^2x-2)}
  {\log_{31}(\sqrt2\cos x)}=0.
  \]

\subsection*{Сложный уровень}

  \item Решите уравнение:
  \[
  \frac{(2\cos x-1)(\sin2x-\sin x)}
  {\sqrt{\sin x}\,\log_{5}(2\cos x)}=0.
  \]
\end{enumerate}

\newpage

\section*{Ответы}

\renewcommand{\arraystretch}{1.45}

\begin{table}[h!]
\centering
\caption{Ответы к тренировочному блоку}
\begin{tabularx}{\linewidth}{@{}cX@{}}
\toprule
\textbf{№} & \textbf{Ответ}\\
\midrule

1 &
\[
x=\pi n,\ n\in\mathbb Z;
\qquad
x=\frac{5\pi}{6}+2\pi k,\ 
x=\frac{7\pi}{6}+2\pi k,\ k\in\mathbb Z.
\]
\\

2 &
\[
x=\frac{\pi}{2}+2\pi k;\qquad
x=\frac{3\pi}{4}+2\pi k;\qquad
x=\frac{5\pi}{4}+2\pi k,\quad k\in\mathbb Z.
\]
\\

3 &
\[
x=\frac{\pi}{2}+\pi k;\qquad
x=\frac{\pi}{6}+2\pi k;\qquad
x=\frac{11\pi}{6}+2\pi k,\quad k\in\mathbb Z.
\]
\\

4 &
\[
x=\frac{\pi}{2}+2\pi k;\qquad
x=\frac{\pi}{3}+2\pi k;\qquad
x=\frac{5\pi}{3}+2\pi k,\quad k\in\mathbb Z.
\]
\\

5 &
\[
x=\frac{\pi}{6}+2\pi k;\qquad
x=\frac{5\pi}{6}+2\pi k,\quad k\in\mathbb Z.
\]
\\

6 &
\[
x=-\frac{\pi}{4}+\pi k,\quad k\in\mathbb Z.
\]
\\

7 &
\[
x=2\pi k;\qquad
x=-\frac{\pi}{3}+2\pi k,\quad k\in\mathbb Z.
\]
\\

8 &
\[
x=\frac{\pi}{3}+2\pi k;\qquad
x=\frac{2\pi}{3}+2\pi k;\qquad
x=\frac{\pi}{2}+2\pi k,\quad k\in\mathbb Z.
\]
\\

9 &
\[
x=-\frac{\pi}{3}+2\pi k,\quad k\in\mathbb Z.
\]
\\

10 &
\[
x=\frac{\pi}{3}+2\pi k,\quad k\in\mathbb Z.
\]
\\

\bottomrule
\end{tabularx}
\end{table}

\end{document}